\chapter{Idea and Concept}
\label{chap:concept}

Chapter~\ref{chap:introduction} formulated the central question: whether a denser point cloud provides more useful evidence to a downstream perception model. Chapter~\ref{chap:fundamentals} defined the geometric and task metrics used to examine this question. The present chapter converts those observations into a controlled experimental design. The design uses two data domains and three downstream networks, but applies the same logical separation between the input visible to the upsampling method, the generated point set, the point set supplied to the downstream model, and the final task metric.

\section{Overall Concept}
\label{sec:fc-3-1}
\label{subsec:metrics_limitations}

The metric definitions explain why a good geometric score is not itself a
guarantee of better perception. CD averages nearest-neighbour discrepancies;
PointNet++ pools learned features within sampled neighbourhoods;
PointRCNN generates proposals from foreground points; and a voxel-based
CenterPoint aggregates points into spatial cells
\cite{qian2021pugcn,shi2019pointrcnn,yin2021centerpoint,qi2017pointnetpp}.
These operations do not optimise or measure the same quantity.

Several consequences follow from the definitions. Repeating coordinates
can increase stored point count without extending spatial coverage. A
small number of misplaced points can have little effect on an average
geometric distance while changing a sampled neighbourhood or an occupied
voxel. Conversely, an improved surface approximation may leave the
predicted class unchanged when the classifier already assigns the correct
class the largest score. These are possible mechanisms derived from the
metrics and architectures, not universal empirical outcomes.

Reference quality also limits interpretation. A complete mesh supports
point-to-surface evaluation, whereas a single sensor scan supplies only
observed returns~\cite{geiger2013kitti,qian2021pugcn}. Distance to that scan
measures consistency with its samples; it cannot certify the unobserved
surface. Geometry, measurement consistency and semantic task performance
must therefore be kept distinct.

For the classification branch, PointNet++ is selected because its set
abstraction layers explicitly sample centroids and aggregate local
neighbourhoods~\cite{qi2017pointnetpp}. Upsampling can change both the
selected centroids and the points grouped around them. It can supply
additional surface support, but near-duplicates need not extend coverage
and displaced points can alter local geometry. This architectural
sensitivity motivates the classifier choice; whether it improves or
degrades classification is determined by the experiments.

The experimental object is the combination of an upsampling method and
its integration protocol. Let $\mathcal X_s^{(b,l)}$ be the visible input
for sample $s$, branch $b\in\{M,K\}$ (ModelNet40 or KITTI), and line
$l\in\{A,B\}$. For method $m$, let $U_m$ be its fixed inference operator;
the method-native candidate set is

\begin{equation}
\mathcal{C}_{s,m}^{(b,l)}=U_m(\mathcal{X}_{s}^{(b,l)}).
\label{eq:fc-3-1}
\end{equation}

The symbol \ensuremath{\mathcal{C}} is used deliberately: before the common output protocol is applied, the method output is treated as a candidate point set. Method-native output cardinality, patch overlap, and representation-specific reconstruction are therefore not automatically equated with the final point cloud evaluated by the downstream task.


The two task branches address complementary settings. ModelNet40 provides
mesh-sampled objects with known categories; KITTI provides partial scene
observations containing foreground objects and background
\cite{wu2015modelnet,geiger2013kitti}. This pairing makes it possible to
examine geometric and classification behaviour within the object branch
before considering transfer to scene-level detection. Agreement between
the branches would support the finding across the tested settings.
Disagreement must be interpreted with their different data preparation
and downstream training regimes in view; it does not identify the data
domain as the sole cause.

\section{Research Hypotheses}
\label{sec:research-hypotheses}

The following hypotheses organise the completed comparisons. They state
relationships to examine rather than conclusions assumed before
evaluation. Their numerical evaluation is presented in
Chapter~\ref{chap:results}.


H1 - Geometric improvement does not imply task improvement. A reduction in CD, HD, P2F, or a related geometric discrepancy is not expected to guarantee an increase in classification accuracy or 3D detection AP.

\begin{samepage}
The implication
\begin{equation}
\Delta d_{\mathrm{geom}}<0\quad\Longrightarrow\quad\Delta\mu_{\mathrm{task}}>0\;?
\label{eq:fc-3-2}
\end{equation}
\end{samepage}

is therefore treated as an empirical question rather than an assumption. Here $d_{\mathrm{geom}}$ denotes a geometric error, $\mu_{\mathrm{task}}$ a task metric such as OA or AP$_{R40}$, and $\Delta$ the method value minus its comparison baseline.

H2 - Densification and recovery are distinct tasks. Adding points to an already available observation and trying to recover information from an artificially sparsified observation are not equivalent questions. A method may be neutral or harmful in the former case but useful in the latter, or vice versa. The two cases are therefore evaluated in separate comparison lines.

H3 - Point count is not equivalent to independent spatial evidence. Satisfying an exact fourfold cardinality requirement only proves that more coordinates are present. It does not prove that those coordinates occupy new, valid surface regions or improve the representation consumed by the downstream model. Duplicate or near-duplicate points can satisfy the count without increasing spatial coverage.

H4 - The downstream architecture mediates the effect of upsampling. PointNet++, PointRCNN, and CenterPoint process point clouds through different sampling and aggregation mechanisms \cite{qi2017pointnetpp,shi2019pointrcnn,yin2021centerpoint}. This motivates the hypothesis that a change helping one architecture may not help another. Cross-model consistency is therefore treated as stronger evidence than an isolated gain in a single architecture.

H5 - Integration and input budgets can dominate the observed effect. Patch locality, normalisation, coordinate restoration, point subsampling, and voxelisation may change the effective dose of generated points. If an upsampled point cloud performs poorly, the experiment must distinguish a failure of the generated geometry from a failure induced by the wrapper or the downstream input adapter.

\section{Evaluation Pipeline}
\label{sec:fc-3-2}

The comparison separates four stages: the information visible to the
upsampler, its generated candidates, the validated output, and the
representation consumed by the task model. Geometry is measured on the
saved output against a declared reference; task performance is measured
after the downstream adapter. Linking these stages by sample identity
allows a geometric change and a task change to be compared without
assuming that every generated point reaches the model.

The two branches address different uses of upsampling. ModelNet40 asks
whether a generated representation supports classification learning;
each representation therefore has its own classifier. The initial KITTI
comparison asks whether an existing detector benefits when upsampling is
inserted before inference, so detector weights remain fixed. A separate
adaptation comparison asks whether training the detector on PU-GCN inputs
reduces the resulting mismatch. It changes the downstream detector while
keeping the upsampler fixed. These training regimes must remain separate
when the findings are combined.

Candidate generation is separated from evaluation. Reference surfaces,
class labels and ground-truth boxes provide assessment targets, not
additional observations for the upsampler. Chapter~\ref{chap:setup}
specifies the transformations, point counts, checkpoints and training
settings that implement this design.

\section{Experimental Comparison}
\label{sec:fc-3-3}

The two comparison lines distinguish adding points to an available
observation from attempting to recover information after controlled
sparsification. They share the same methods within each branch, but
their baselines answer different questions. Table~\ref{tab:lines} in
Chapter~\ref{chap:setup} specifies their executed point counts.

\subsection{Line A - Densification of the Original Observation}
\label{sec:fc-3-3-1}

Line~A asks whether generated points provide useful marginal information
when the original observation is already available. Its baseline is the
native original input, without an extra resampling operation. The
comparison therefore concerns the practical consequence of adding the
upsampling stage, rather than a comparison between two artificially
equalised versions of the baseline.

A negative Line~A result cannot be attributed to deliberately removing
observations beforehand. It means that the tested generator, integration
and downstream model did not obtain a net task benefit from densifying
the available input. It does not imply that every individual object or
category deteriorates.

\subsection{Line B - Recovery after Controlled Sparsification}
\label{sec:fc-3-3-2}

Line~B removes a fixed proportion of observations before upsampling.
Every method receives the same prepared sparse input. Its immediate
baseline is that sparse representation, so a gain measures recovery
relative to the information actually available to the generator. The
original representation remains a separate pre-sparsification reference.

This reduction is a controlled stress condition, not a simulation of
a particular lower-channel LiDAR sensor. Restoring the number of points
does not restore the missing measurements themselves. The comparison
therefore asks how much useful task information a method recovers,
rather than assuming that an equal output count establishes recovery.

\section{Fairness Protocol}
\label{sec:fc-3-4}
\label{sec:fc-3-4-1}
\label{sec:fc-3-4-2}
\label{sec:fc-3-4-3}
\label{sec:upsampling-inference-fairness}
\label{sec:fc-3-4-4}
\label{sec:fc-3-6}

Fairness requires the same visible observations, a declared output
cardinality and a consistent evaluation within each comparison.
Upsampling methods share the target expansion ratio; native baselines
retain their original sizes. Geometry comparisons separately control
the evaluated and reference counts. This distinguishes equal conditions
among generated representations from an artificial enlargement of the
baseline.

A common interface does not imply identical internal processing.
PU-Net expands learned point features, PU-GCN uses graph features and
NodeShuffle, PU-EdgeFormer uses edge-transformer features, PDANS uses
diffusion-based generation, and the EAR-style baseline uses geometric
resampling
\cite{yu2018punet,qian2021pugcn,kim2023puedgeformer,zhang2025pdans,huang2013ear}.
Their implementation-specific preprocessing must therefore be recorded
alongside the common comparison rules. A modified inference path supports
a conclusion about that integrated pipeline, rather than automatically
about an unmodified published method.

The effective input depends on the downstream architecture. PointRCNN
samples points, CenterPoint aggregates them into voxels, and PointNet++
groups local neighbourhoods
\cite{shi2019pointrcnn,yin2021centerpoint,qi2017pointnetpp}. Let
$\mathcal Y_{s,m}$ denote the validated output for sample $s$ and
upsampler $m$. The representation consumed by downstream model $d$ is
\begin{equation}
\widetilde{\mathcal Y}^{(d)}_{s,m}=A_d(\mathcal Y_{s,m}),
\label{eq:fc-3-8}
\end{equation}
where $A_d$ is its input adapter. Equal saved cardinality need not imply
equal distinct support after this operation. Adapter settings and
diagnostic changes are consequently reported as part of the experimental
protocol, rather than inferred from file size.

Evidence is interpreted at its actual scope. Complete dataset comparisons
establish the main task outcomes; detector adaptation tests a separate
training intervention; targeted diagnostics examine specific integration
factors. Subset results do not replace full-validation results, and an
adapted input must be compared with the corresponding adapted control
where one exists. Chapter~\ref{chap:setup} records which controls were
completed and which sources of uncertainty remain.

\section{Cross-Domain Interpretation}
\label{sec:fc-3-5}
\label{sec:fc-3-7}

Object surfaces and LiDAR scenes differ in visibility, coverage and
sampling. Scene occlusion and range-dependent support
\cite{geiger2013kitti} make the transfer of object-oriented upsamplers
an additional experimental question
\cite{yu2018punet,qian2021pugcn,kim2023puedgeformer,zhang2025pdans}.
Patch locality, restoration of metric coordinates and preservation of
measured attributes can mediate that transfer. Their implementations
belong to Chapter~\ref{chap:setup}; the conceptual requirement is to
distinguish these interface effects from the generated geometry itself.

Three comparisons connect the branches. Within each branch, Line~A and
Line~B distinguish densification from recovery. Within ModelNet40,
geometric ordering is compared with classification ordering. Across
branches, the direction of the task response is compared while retaining
the different training regimes and observation types. Overall accuracy
and detection AP are not combined into a universal score.

Downstream reception and voxel occupancy provide intermediate
measurements that may explain a task response. They do not by themselves
identify a unique cause, and selected visual cases do not estimate how
frequently an outcome occurs. The design thus separates available
information, generated geometry, output assembly and downstream use.
The next chapter specifies how those distinctions were implemented in
the completed experiments.
